% Generated draft: requires mathematical and source review.
\documentclass[11pt]{article}
\usepackage[margin=25mm]{geometry}
\usepackage{fontspec}
\setmainfont{Latin Modern Roman}
\usepackage{amsmath,amssymb,mathtools,unicode-math}
\setmathfont{Latin Modern Math}
\usepackage{enumitem,needspace,xurl,hyperref}
\setlength{\parindent}{0pt}
\setlength{\parskip}{5pt}
\setlength{\emergencystretch}{3em}
\newcommand{\N}{\mathbb{N}}
\newcommand{\Z}{\mathbb{Z}}
\newcommand{\Q}{\mathbb{Q}}
\newcommand{\R}{\mathbb{R}}
\newcommand{\C}{\mathbb{C}}
\newcommand{\F}{\mathbb{F}}
\newcommand{\A}{\mathbb{A}}
\newcommand{\PP}{\mathbb{P}}
\newcommand{\E}{\mathbb{E}}
\newcommand{\eps}{\varepsilon}
\newcommand{\dd}{\,\mathrm{d}}
\begin{document}
\section*{UnsolvedMath ID: 801012}
\textbf{Source problem code:} catalogue entry
\subsection{Definitions and mathematical statement}
\textbf{Interpretive qualifications.} The supplied record does not explicitly specify the domain of the sequences or define the complexity and reducibility notation. For a conditional standard interpretation, take \(x,z\in 2^{\N}\), the set of infinite binary sequences, and let \(x\upharpoonright n\) denote the initial segment of \(x\) of length \(n\). Let \(K(\sigma)\) denote prefix-free Kolmogorov complexity of a finite binary string \(\sigma\), relative to a fixed optimal prefix-free machine, and let \(K(n)\) denote the complexity of a fixed effective encoding of the natural number \(n\). The usual initial-segment complexity ordering is
\[
z\leq_K x
\quad\Longleftrightarrow\quad
(\exists c\in\N)(\forall n\geq 1)\;
K(z\upharpoonright n)\leq K(x\upharpoonright n)+c,
\]
with \(z\geq_K x\) meaning \(x\leq_K z\). These are explanatory conventions, whose agreement with the intended source definitions requires review.

The source uses the distinct notation \(\leq_K^+\) in its lower cone. Its superscript \(+\) is not explained in the supplied record; no identification with \(\leq_K\) is asserted here. Retaining that unresolved notation, set
\[
L_x=\{z\in 2^{\N}:z\leq_K^+x\},
\qquad
U_x=\{z\in 2^{\N}:z\geq_Kx\}.
\]
Here countable includes finite sets.

\textbf{Questions, subject to the preceding interpretive qualifications.} Are both of the following assertions true for every \(x\in 2^{\N}\)?
\begin{enumerate}
\item
\[
L_x\text{ is countable}
\quad\Longleftrightarrow\quad
\liminf_{n\to\infty}\bigl(K(x\upharpoonright n)-K(n)\bigr)<\infty.
\]
\item
\[
U_x\text{ is countable}
\quad\Longleftrightarrow\quad
\liminf_{n\to\infty}\bigl(K(n)+n-K(x\upharpoonright n)\bigr)<\infty.
\]
\end{enumerate}
Each right-hand condition asks for a uniform finite upper bound along infinitely many lengths. The intended domain and, especially, the meaning of \(\leq_K^+\) must be confirmed before this can be treated as a fully specified formulation of the source problem.
\subsection{Short English statement}
The problem asks whether countability of the lower and upper complexity cones of an infinite sequence is characterized by the two displayed bounds holding infinitely often. The binary-sequence and prefix-free-complexity conventions are conditional, and the source's notation \(\leq_K^+\) remains undefined. The registry labels the problem open; that label is supplied metadata, not an independently verified assessment.
\subsection{Sources}
\begin{itemize}
\item Local UnsolvedMath registry: supplied complete source record, permanent ID \(801012\), exact source problem code \texttt{catalogue entry}; title: \emph{Open Problems in Mathematical Logic --- Problem 24}. The record attributes the problem to Barmpalias, Gao et al., with proposed year \(2026\), and supplies the status \texttt{open}.
\item Barmpalias, Gao et al., \emph{Open Problems in Mathematical Logic} (2026), Problem 24. Supplied source URL: \url{https://arxiv.org/abs/2608.26628}. The record reports access on \texttt{2026-10-16} and extraction by \texttt{arxiv\_latex}. These are supplied provenance claims; the source was not independently accessed or verified for this draft.
\end{itemize}
\end{document}
